% =====================================================================
% SEÇÃO 4 E 5 — DESAFIOS TECNOLÓGICOS E CONCLUSÃO
% Estilo de Alto Rigor Acadêmico (Padrão IEEE / Gabriel Vidal Gaspar)
% =====================================================================

\section{Desafios Tecnológicos e Limitações do Sistema}
\label{sec:desafios}

Seguindo a perspectiva crítica de tradução tecnológica estabelecida na literatura recente, o pipeline agnóstico BioStatusIA apresenta limitações estruturais que delimitam sua aplicação atual:
\begin{enumerate}
    \item \textbf{Gargalo Computacional em Volumes 3D (F4):} O cálculo de descritores radiômicos tridimensionais (GLCM em 3 planos ortogonais e esfericidade 3D) impõe alto consumo de memória RAM e GPU quando executado sobre tomografias/ressonâncias de altíssima resolução.
    \item \textbf{Fronteira de Escopo de Modalidades:} Modalidades complexas não contempladas na versão atual (como sequenciamento genômico ou vídeo ultrassom contínuo) exigem a implementação de leitores e extratores dedicados.
    \item \textbf{Trade-off entre Flexibilidade e Especialização:} Embora o pipeline unificado elimine o custo de engenharia de software para suporte a novas modalidades, pipelines ultra-especializados e customizados manualmente para uma única tarefa clínica podem obter ganhos marginais de desempenho.
\end{enumerate}

\section{Conclusão e Trabalhos Futuros}
\label{sec:conclusao}

Este trabalho apresentou o projeto, a implementação e a avaliação empírica de um pipeline de ingestão e engenharia de features biomédicas agnóstico à modalidade. A arquitetura demonstrou robustez ideal ($100\%$ de sucesso sem exceções não tratadas) no estresse experimental com **30 bases de dados biomédicas de benchmark reais**, abrangendo sinais fisiológicos temporais (F1), DICOM 2D (F3), volumes 3D (F4), imagens convencionais 2D e dados tabulares.

A introdução do algoritmo determinístico de 9 modos de detecção estrutural e do contrato unificado \texttt{SinalNormalizado} comprovou que é possível desacoplar os leitores físicos dos algoritmos analíticos de extração de biomarcadores. Além disso, os resultados demonstraram que a dimensionalidade do vetor de features gerado é uma função invariante da família biomédica.

Como trabalhos futuros, esta pesquisa fundamenta o desenvolvimento do **Paper B**, focado no treinamento de modelos AutoML (classificação, seleção de hiperparâmetros, calibração de probabilidades e interpretabilidade SHAP) e na geração automatizada de laudos clínicos por inteligência artificial a partir do vetor de features produzido nesta etapa.
